\chapter{Komposition und gemeinsame Zuordnungen}
\label{b05guide:ch:composition}

\section{Hintereinanderausführen von Funktionen}

\begin{definition}[Komposition]
\label{b05guide:def:composition}
Für $f\colon A\to B$ und $g\colon B\to C$ heißt
\[
 g\circ f
   =\{(a,c)\in A\times C\mid
       \exists b\in B\,((a,b)\in f\land(b,c)\in g)\}
\]
die \emph{Komposition} von $f$ und $g$.
Dabei wird zuerst $f$, danach $g$ ausgeführt.
\end{definition}

\begin{theorem}[Typisierung und Auswertung der Komposition]
\label{b05guide:thm:composition-value}
Für $f\colon A\to B$ und $g\colon B\to C$ gilt
\[
 g\circ f\colon A\to C,\qquad
 \forall a\in A\,((g\circ f)(a)=g(f(a))).
\]
Die Komposition ist die eindeutig bestimmte Funktion auf $A$
mit diesen Werten. Die Aussage gilt auch für leere Mengen.
\end{theorem}

\begin{theorem}[Assoziativität]
\label{b05guide:thm:composition-associativity}
Für $f\colon A\to B$, $g\colon B\to C$ und
$h\colon C\to D$ gilt
\[
 h\circ(g\circ f)=(h\circ g)\circ f\colon A\to D.
\]
An jeder Stelle $a\in A$ haben beide Graphen den Wert
$h(g(f(a)))$.
\end{theorem}

\begin{theorem}[Bilder und Urbilder einer Komposition]
\label{b05guide:thm:composition-images}
Seien $f\colon A\to B$, $g\colon B\to C$,
$U\subseteq A$ und $W\subseteq C$. Dann gilt
\[
 \begin{aligned}
 (g\circ f)[U]&=g[f[U]],\\
 (g\circ f)^{-1}[W]&=f^{-1}[g^{-1}[W]].
 \end{aligned}
\]
Insbesondere ist $(g\circ f)[A]=g[f[A]]$.
\end{theorem}

\begin{theorem}[Einschränkungen und Zwischenmengen]
\label{b05guide:thm:composition-restrictions}
Seien $f\colon A\to B$, $g\colon B\to C$ und
$U\subseteq A$. Dann gilt
\[
 (g\circ f)\restriction_U=g\circ(f\restriction_U).
\]
Eine Teilmenge $D\subseteq B$ ist für dieselbe Funktion $f$
genau dann als Zwischenmenge zulässig, wenn $f[A]\subseteq D$:
\[
 f\colon A\to D\quad\Longleftrightarrow\quad f[A]\subseteq D.
\]
In diesem Fall gilt, mit $f$ als Funktion nach $D$,
\[
 (g\restriction_D)\circ f=g\circ f\colon A\to C.
\]
Für jedes $D\subseteq B$ gilt ohne weitere Voraussetzung
\[
 (g\restriction_D)\circ(f\restriction_{f^{-1}[D]})
       =(g\circ f)\restriction_{f^{-1}[D]},
\]
wobei $f\restriction_{f^{-1}[D]}$ als Funktion nach $D$
aufgefaßt wird. Die gemeinsame Ausgangsmenge ist $f^{-1}[D]$.
\end{theorem}

\begin{theorem}[Gleichheit nach einem vorgeschalteten Schritt]
\label{b05guide:thm:composition-equality}
Für $f\colon A\to B$ und $g,h\colon B\to C$ gilt
\[
 \begin{aligned}
 g\circ f=h\circ f
 &\quad\Longleftrightarrow\quad
       \forall b\in f[A]\,(g(b)=h(b))\\
 &\quad\Longleftrightarrow\quad
       g\restriction_{f[A]}=h\restriction_{f[A]}.
 \end{aligned}
\]
Die Werte von $g$ und $h$ außerhalb von $f[A]$ gehen
in diese Kompositionen nicht ein. Für $A=\varnothing$
sind alle drei Bedingungen erfüllt.
\end{theorem}

\section{Zwei Werte gemeinsam zuordnen}

\begin{definition}[Gepaarte Funktion]
\label{b05guide:def:paired-function}
Für $f\colon A\to B$ und $g\colon A\to C$ sei
\[
 \langle f,g\rangle
   =\{(a,(f(a),g(a)))\mid a\in A\}.
\]
Diese Konstruktion faßt die beiden Werte zu einem
geordneten Paar zusammen.
\end{definition}

\begin{theorem}[Existenz und Eindeutigkeit der gepaarten Funktion]
\label{b05guide:thm:paired-function}
Für $f\colon A\to B$ und $g\colon A\to C$ ist
\[
 \langle f,g\rangle\colon A\to B\times C,\qquad
 \langle f,g\rangle(a)=(f(a),g(a))\quad(a\in A).
\]
Es ist die eindeutig bestimmte Funktion mit dieser
Wertebeschreibung. Für $U\subseteq A$ gilt
\[
 \langle f,g\rangle[U]
   =\{(f(a),g(a))\mid a\in U\}
   \subseteq f[U]\times g[U].
\]
Für $V\subseteq B$ und $W\subseteq C$ gilt
\[
 \langle f,g\rangle^{-1}[V\times W]
      =f^{-1}[V]\cap g^{-1}[W].
\]
\end{theorem}

\begin{definition}[Produktfunktion]
\label{b05guide:def:product-function}
Für $f\colon A\to B$ und $g\colon C\to D$ sei
\[
 f\boxtimes g
  =\{((a,c),(f(a),g(c)))\mid a\in A,\ c\in C\}.
\]
Das Zeichen $\boxtimes$ bezeichnet diesen Funktionsgraphen,
nicht das kartesische Produkt der beiden Graphen.
\end{definition}

\begin{theorem}[Werte, Bilder und Urbilder der Produktfunktion]
\label{b05guide:thm:product-function}
Für $f\colon A\to B$ und $g\colon C\to D$ gilt
\[
 \begin{gathered}
 f\boxtimes g\colon A\times C\to B\times D,\\
 (f\boxtimes g)(a,c)=(f(a),g(c))
       \qquad(a\in A,\ c\in C).
 \end{gathered}
\]
Die Funktion ist durch diese Werte eindeutig bestimmt.
Für $U\subseteq A$, $V\subseteq C$, $S\subseteq B$
und $T\subseteq D$ gilt
\[
 \begin{aligned}
 (f\boxtimes g)[U\times V]&=f[U]\times g[V],\\
 (f\boxtimes g)^{-1}[S\times T]
     &=f^{-1}[S]\times g^{-1}[T].
 \end{aligned}
\]
Ist $A=\varnothing$ oder $C=\varnothing$, so ist
$f\boxtimes g$ die leere Funktion auf $A\times C=\varnothing$.
\end{theorem}

\begin{theorem}[Verträglichkeit mit Komposition]
\label{b05guide:thm:pair-product-composition}
Für $h\colon X\to A$, $f\colon A\to B$ und $g\colon A\to C$
gilt
\[
 \langle f,g\rangle\circ h
     =\langle f\circ h,g\circ h\rangle.
\]
Für $f\colon A\to B$, $g\colon C\to D$,
$u\colon B\to E$ und $v\colon D\to F$ gilt
\[
 (u\boxtimes v)\circ(f\boxtimes g)
       =(u\circ f)\boxtimes(v\circ g).
\]
Beide Aussagen gelten ohne zusätzliche Eigenschaften
der beteiligten Funktionen.
\end{theorem}
